\documentclass[11pt]{article}
\usepackage[T1]{fontenc}
\usepackage{lmodern,mathrsfs,amsmath,amssymb,amsthm,mathtools,booktabs,microtype}
\usepackage[margin=1in]{geometry}
\usepackage[hidelinks]{hyperref}
\setlength{\emergencystretch}{3em}
\newtheorem{theorem}{Theorem}[section]
\newtheorem{lemma}[theorem]{Lemma}
\newtheorem{proposition}[theorem]{Proposition}
\newtheorem{corollary}[theorem]{Corollary}
\theoremstyle{definition}\newtheorem{definition}[theorem]{Definition}
\theoremstyle{remark}\newtheorem{remark}[theorem]{Remark}
\newcommand{\R}{\mathbb R}\newcommand{\Q}{\mathbb Q}\newcommand{\Z}{\mathbb Z}
\newcommand{\CP}{\mathcal{CP}}\newcommand{\COP}{\mathcal{COP}}
\newcommand{\SPN}{\mathcal{SPN}}\newcommand{\Ry}{\mathcal R}
\newcommand{\tr}{\operatorname{tr}}\newcommand{\diag}{\operatorname{diag}}
\newcommand{\cone}{\operatorname{cone}}\newcommand{\Min}{\operatorname{Min}}
\newcommand{\ip}[2]{\langle #1,#2\rangle}\newcommand{\norm}[1]{\lVert #1\rVert}
\newcommand{\trans}{^{\mathsf T}}
\title{Long pivot paths and short certificates for completely positive matrices\\[3pt]\large Finite local models, a normal-cone lifting theorem, and intrinsic graph growth}
\author{Manuscript II research revision}
\date{15 November 2026}
\begin{document}\maketitle
\begin{abstract}
We give a verifiable sufficient-conditions theorem converting a completely positive nonmember into exponentially long improving paths on the copositive Ryshkov polyhedron. The leading block of the seed must expose a fixed positive-definite perfect form; no particular graph or exceptional copositive matrix is required. The geometric mechanism is a finite-template theorem: after an integral shear, every active vector and every primitive tangent ray along a fixed-block face belongs to a bounded finite set. Consequently a scale linear in the prescribed path length suffices, and each individual local step has polynomial bit complexity in fixed dimension. The bound covers every strictly improving adjacent rule and any fixed finite menu of starts on the face. It survives arbitrary positive-semidefinite corrections of local rays, by a classical copositive decomposition lemma. New qualifying families exist in every odd order and have all proper principal submatrices completely positive, excluding all lower-order principal certificate tests. Independently, every fixed-positive-definite-block quadratic epigraph has bounded degree and intrinsic graph-ball growth of order $R^{n-1}$, with exponentially separated vertices of logarithmic encoding length. The original canonical five-cycle construction is recovered without changing its validity.
\end{abstract}

\section{Introduction and the lifting principle}
The ordinary completely positive and copositive cones are
\[
 \CP_n=\cone\{xx\trans:x\ge0\},\qquad
 \COP_n=\{P:x\trans Px\ge0\ (x\ge0)\}.
\]
Their pairing is $\ip A P=\tr(AP)$. Let
\[
 \Ry_n=\{P:P[v]\ge1\text{ for }0\ne v\in\Z_+^n\},\qquad P[v]=v\trans Pv.
\]
A vertex is a perfect copositive form. Its finitely many active vectors generate its minimizing normal cone. These objects underlie the simplex algorithm of Dutour Sikiri\'c--Sch\"urmann--Vallentin \cite{DSV}; their geometry is developed in \cite{DS,OS}. We study the local model that chooses an arbitrary extreme tangent ray of strictly negative objective cost and traverses its bounded edge. We do not restrict the computation used to choose the ray.

The preceding version of this manuscript treated a specific five-cycle seed from the canonical tridiagonal vertex. It used bounded edge movement and location-dependent cofactor estimates. Here the main object is instead the entire fixed-block face. Its finite local models yield a stronger scale bound, a criterion for arbitrary seeds, and an independent graph-growth theorem. Odd cycles are applications of the criterion, not assumptions of it.

Put $r=n-1$, $d=n(n+1)/2$. Fix a rational positive-definite perfect form $H\in\Ry_r$, and define
\[
 \mathcal M_H=\{v\in\Z_+^r:H[v]=1\},\qquad
 \mathcal V_H=\cone\{vv\trans:v\in\mathcal M_H\}.
\]
The set $\mathcal M_H$ is finite and its squares span $\mathbb S^r$. The face obtained by keeping these equalities in order $n$ is
\begin{equation}\label{eq:face}
 \mathcal F_H=\left\{P(b,c)=\begin{pmatrix}H&b\\b\trans&c\end{pmatrix}\in\Ry_n\right\}.
\end{equation}
It is nonempty and has dimension $r+1$; Section~\ref{sec:geometry} proves the needed local properties.

\begin{definition}[Qualifying seed]\label{def:seed}
A qualifying seed for $H$ is a rational matrix
\[
 B=\begin{pmatrix}B_0&a\\a\trans&\alpha\end{pmatrix}
\]
which is positive definite, strictly entrywise positive, outside $\CP_n$, and satisfies $B_0\in\operatorname{int}\mathcal V_H$.
\end{definition}
This is a finite verifiable criterion in each fixed dimension. Positive definiteness is an exact rational test. Strict exposure is a finite LP: require a representation of $B_0$ with positive coefficients on every vector in $\mathcal M_H$. Nonmembership can be supplied with a rational $W\in\COP_n$ and $\ip W B<0$. Copositivity of a rational seed witness can be checked by minimizing on all simplex supports; no variable-dimensional polynomial claim is made. The criterion does not require a five-cycle witness.

Here are explicit constants for the main theorem. Choose an integer $q_0\ge2$ with $q_0^2\ge2r\norm H_2$, and rational upper bounds
\[
 A_H\ge\max_i\sqrt{(H^{-1})_{ii}},\quad
 L_H\ge\max_i(1/2+2\sqrt{H_{ii}}).
\]
Set
\begin{equation}\label{eq:constants}
 \begin{split}
 Y&=2q_0^r,\\
 M_H&=\left\lceil\max\{1,Y+A_H,\max_{v\in\mathcal M_H}\norm v_\infty\}\right\rceil,\\
 F_H&=(d-1)![2M_H^2]^{d-1},\\
 K_H&=\norm{H^{-1}}_\infty.
 \end{split}
\end{equation}
The matrix norm in $K_H$ is the maximum absolute row sum. Let $q_B$ be a positive integer clearing the denominators of $B_0$, and let $R_0$ bound $\norm b_\infty$ at every permitted initial vertex. Put
\begin{equation}\label{eq:scale}
 \begin{split}
 T_B&=2\norm a_1+\alpha+n+1,\\
 s_m&=1+rK_H(R_0+mL_H+1/2),\\
 N_m&=\left\lceil8q_BF_HT_Bs_m\right\rceil,\\
 S_N&=\diag(NI_r,1).
 \end{split}
\end{equation}
In particular $N_m=\Theta_{H,B,R_0}(m+1)$.

For robustness, fix rational data
\[
 B\succeq\mu I,\quad b_*\le\min_{i,j}B_{ij},\quad
 W\in\COP_n,\quad\ip W B\le-\gamma,\quad\norm W_F\le\omega,
\]
where $\mu,b_*,\gamma>0$. Let $\beta\ge\max(1,\norm B_F)$ be rational, and define
\begin{equation}\label{eq:robust}
 t_B=\frac\gamma{4\tr B},\quad w_B=\omega+t_B,\quad
 \eta=\min\{1,\mu/2,b_*/2,\gamma/(4w_B)\},\quad
 A_N=\frac{S_NBS_N}{2\beta N^2}.
\end{equation}
Shrinking $\eta$ is always allowed.

\begin{theorem}[Normal-cone lifting; all local choices]\label{thm:main}
Let $H,B$ satisfy Definition~\ref{def:seed}, and let the initial vertex be any vertex of $\mathcal F_H$ with $\norm b_\infty\le R_0$. For every $m\ge1$, $N\ge N_m$, and symmetric $A$ with
\begin{equation}\label{eq:ball}
 \norm{A-A_N}_F\le\frac\eta{2\beta N^2},
\end{equation}
the following hold.
\begin{enumerate}
\item $A$ is positive definite, strictly entrywise positive, has norm at most one, and is not completely positive.
\item Every strictly improving extreme-ray adjacent rule makes at least $m$ pivots. Throughout this prefix every improving ray is tangent to $\mathcal F_H$ and bounded; an improving ray exists at each departure vertex.
\item No copositive matrix whose leading $r\times r$ principal block is PSD separates $A$. This excludes arbitrary PSD corrections, hence arbitrary nonnegative diagonal corrections, of any visited improving ray. It also excludes history-based certificates
\[
 \sum_j\lambda_jP_j+\sum_k\theta_kR_k+Q,
 \qquad \sum_j\lambda_j\ge0,\quad\theta_k\in\R,\quad Q\succeq0,
\]
whenever they are copositive. The $P_j$ may come from different runs on the same face.
\item The rational matrix
\[
 W_N=S_N^{-1}WS_N^{-1}+\frac{t_B}{N^2}I
\]
works for the whole ball, with
\[
 W_N[x]\ge\frac{t_B}{N^2}\norm x_2^2\ (x\ge0),\quad
 \norm{W_N}_F\le w_B,\quad
 \ip{W_N}A\le-\frac\gamma{4\beta N^2}.
\]
In particular the distance from $A$ to $\CP_n$ is at least $\gamma/(4\beta w_BN^2)$.
\end{enumerate}
For fixed seed and block, the centers $A_{N_m}$ have binary input-and-distance precision $\Theta(\log(m+1))$, while the forced path has length at least $m$. Each step inside the face, including the full active set, tangent rays, and the selected neighbor, can be computed in polynomial bit time in fixed dimension. A fixed finite menu of starts, parallel runs, and restarts within the stated initial region do not remove the required prefix on any successful run.
\end{theorem}
The number of starts need not be fixed in the assertion about the region; there are only finitely many vertices in that region. The theorem does not cover menus using different leading blocks, nonextreme descent directions, or unrestricted nonlocal separator searches. Its explicit separator is a global bypass. The robust encoding bound concerns the rational centers, not arbitrary rational perturbations.

\section{Finite local models for a fixed quadratic block}\label{sec:geometry}
The geometric statements in this section require only $H\succ0$, not perfection or an objective. Define
\[
 g_H(b)=\sup_{u\in\Z_+^r,\ y\in\Z_{\ge1}}
 \frac{1-u\trans Hu-2yb\trans u}{y^2},\qquad
 \mathcal E_H=\{(b,c):c\ge g_H(b)\}.
\]
Let $z$ minimize $x\trans Hx+2b\trans x$ over $x\ge0$, and write $\ell=b+Hz$. Strong convexity gives a unique minimizer and
\begin{equation}\label{eq:kkt}
 z,\ell\ge0,\qquad z_i\ell_i=0.
\end{equation}
Take $q_0,Y,A_H$ as above and $\kappa=(2q_0^{2r})^{-1}$.

\begin{lemma}[Transverse gap, active denominators, and edge movement]\label{lem:gap}
The supremum defining $g_H$ is attained and
\[
 z\trans Hz+\kappa\le g_H(b)\le z\trans Hz+1.
\]
At an active vector $(u,y)$, writing $e=u-yz$, one has
\begin{equation}\label{eq:active}
 y\le Y,\qquad e\trans He+2y\ell\trans u
 =1-y^2(g_H(b)-z\trans Hz)\le1,\qquad |e_i|\le A_H.
\end{equation}
The epigraph is locally polyhedral. At a vertex, $0\le\ell_i\le1/2$ for every $i$. Along every bounded edge,
\begin{equation}\label{eq:movement}
 |b_i-b_i'|\le1/2+2\sqrt{H_{ii}}.
\end{equation}
Its recession cone is $\{(w,s):w\ge0,\ s\ge0\}$.
\end{lemma}
\begin{proof}
Completing the square gives
\[
 \frac{1-H[u]-2yb\trans u}{y^2}
 =H[z]+\frac{1-H[e]-2y\ell\trans u}{y^2}.
\]
This proves the upper bound. On the positive support of $z$, put the $q_0^s+1$ fractional vectors $jz$, $0\le j\le q_0^s$, into $q_0^s$ cubes of side $1/q_0$. Subtract two in the same cube. This gives $1\le y\le q_0^r$ and integers $u_i$ with $|u_i-yz_i|\le1/q_0$ on that support. The approximants are nonnegative; set the other coordinates to zero. Thus $\ell\trans u=0$ and $H[e]\le1/2$. If $z=0$, use $u=0,y=1$. This proves the lower bound. Denominators larger than $Y$ cannot attain the maximum. The finitely many remaining denominators have possible maximizing numerators in a bounded ellipsoid around $yz$. This proves attainment, \eqref{eq:active}, and local polyhedrality; the minimizer $z(b)$ is continuous by strong convexity.

At a vertex the active normals $(2u/y,1)$ span $\R^{r+1}$. For each coordinate $i$, an active $u_i$ is positive; otherwise its cross-column variable is unconstrained in the active system. Equation~\eqref{eq:active} yields $2y\ell_i u_i\le1$ and hence $\ell_i\le1/2$.

Endpoints of a bounded edge share an active $(u,y)$. At either endpoint,
\[
 b+Hu/y=\ell+He/y,
\]
and $|(He)_i|\le\sqrt{H_{ii}}$ by Cauchy--Schwarz in the $H$-metric. Both values of the left side lie in the same interval
$[-\sqrt{H_{ii}}/y,\,1/2+\sqrt{H_{ii}}/y]$. This proves \eqref{eq:movement}. Finally a recession direction satisfies $s+2w\trans u/y\ge0$ for all $u,y$. Taking $u=0$, and then $u=je_i,y=1$ with $j\to\infty$, proves the asserted signs; they are sufficient.
\end{proof}

For an integer vector $k$, define the shear
\begin{equation}\label{eq:shear}
 \mathsf S_k(b,c)=(b-Hk,\ c-2k\trans b+k\trans Hk),\qquad
 U_k=\begin{pmatrix}I_r&-k\\0&1\end{pmatrix}.
\end{equation}
It acts on matrices by $P\mapsto U_k\trans PU_k$. It need not preserve the whole orthant Ryshkov set. We use it to normalize local data, not as an unqualified symmetry.

\begin{theorem}[Finite-template theorem]\label{thm:templates}
Fix $H\succ0$. Every vertex $(b,c)$ of $\mathcal E_H$ can be written
\[
 (b,c)=\mathsf S_k(\bar b,\bar c),\qquad k\in\Z_+^r,
\]
where $(\bar b,\bar c)$ belongs to a finite set depending only on $H$. After this shear every active vector has integer coordinates bounded in absolute value by $Y+A_H$ and has last coordinate in $\{1,\ldots,Y\}$. The graph of bounded edges has uniformly bounded degree. If $H$ is rational, all vertices have a common bounded denominator.

If additionally $H$ is perfect in $\Ry_r$, the \emph{full} tangent cones of $\Ry_n$ at vertices of $\mathcal F_H$ have finitely many local models under these integral shears. Each primitive integral extreme tangent ray has the form
\begin{equation}\label{eq:raymodel}
 R=U_k\trans\bar R U_k,\qquad
 \bar R=\begin{pmatrix}X&w\\w\trans&s\end{pmatrix},\qquad
 |\bar R_{ij}|\le F_H.
\end{equation}
Consequently, with $v_k=1+r\norm k_\infty$,
\begin{equation}\label{eq:raybounds}
 |X_{ij}|\le F_H,\quad |R_{in}|\le F_Hv_k,\quad
 |R_{nn}|\le F_Hv_k^2,\quad \norm R_F\le nF_Hv_k^2.
\end{equation}
In particular the number of active vectors and the number of competing rays are bounded independently of position on the face.
\end{theorem}
\begin{proof}
At a vertex put $k_i=\lfloor z_i\rfloor$ and $\bar z=z-k\in[0,1)^r$. Set
\[
 \bar b=b+Hk,\qquad \bar c=c+2k\trans b+k\trans Hk.
\]
Complementarity gives $k\trans\ell=0$, so
\[
 \bar b=\ell-H\bar z,\qquad
 \bar c=H[\bar z]+c-H[z].
\]
These coordinates lie in a fixed compact set. An active vector transforms to $(\bar u,y)=(u-yk,y)$, and
$|\bar u_i|\le y+A_H\le Y+A_H$. There are finitely many such integer vectors. At a vertex $r+1$ independent active affine equations determine $\bar b,\bar c$; therefore only finitely many normalized points occur. Some of them need not be feasible orthant vertices before the shear, which does not affect the conclusion. Local tangent cones are specified by subsets of this finite list, giving a uniform degree bound. For rational $H$, the finitely many normalized solutions are rational, and integral shears preserve a common denominator.

For a full matrix vertex, include the fixed active vectors $(v,0)$, $v\in\mathcal M_H$. This finite set together with the transformed positive-last-coordinate vectors lies in the box of radius $M_H$ from \eqref{eq:constants}. The transformed tangent cone is full dimensional and pointed. An extreme ray annihilates $d-1$ independent rank-one equations whose integer coefficients are bounded by $2M_H^2$. Cofactors yield a primitive integral representative bounded by $F_H$. Congruence by the integral unimodular $U_k$ and its integral inverse preserves primitivity. Expanding the blocks of $U_k\trans\bar R U_k$ proves \eqref{eq:raybounds}.
\end{proof}

For clarity about the face used above, the active equations of $H$ force the leading block in \eqref{eq:face} to equal $H$, and the remaining inequalities give exactly $\mathcal E_H$. Every point of this face is strictly copositive: for $a=c-H[z]\ge\kappa$,
\[
 P[(x,y)]=H[x-yz]+2y\ell\trans x+ay^2>0\quad(0\ne(x,y)\ge0).
\]
A neighborhood therefore has a common strict margin; only finitely many integer constraints matter there. Vertices and edges of the face are genuine vertices and edges of the full Ryshkov set. Finite edge endpoints cannot disappear into a nonattained boundary.

\begin{corollary}[Cheap complete local information]\label{cor:localcost}
For fixed rational $H$, the complete active set, full tangent rays at a vertex of $\mathcal F_H$, and the neighbor along any bounded face edge can be computed in polynomial bit time in the encoding of the vertex and ray. The counts of lattice tests and linear systems depend on $H,n$, not on the vertex height.
\end{corollary}
\begin{proof}
Find $z$ by solving the $2^r$ possible support systems of the positive-definite quadratic program and checking \eqref{eq:kkt}. This uses rational arithmetic. Compute $k=\lfloor z\rfloor$ and enumerate the finite normalized box of Theorem~\ref{thm:templates}, retaining $(\bar u+yk,y)\ge0$ and testing equality. Enumerate the extreme rays from the finite active system. For a bounded face neighbor, \eqref{eq:movement} and $z=H^{-1}(\ell-b)$ bound the change of $z$ by $K_H(L_H+1/2)$ in sup norm. Thus the endpoint's active numerators, normalized by the \emph{old} $k$, lie in another fixed box. Enumerate their positive constraint ratios along the ray. The smallest is the exact edge length because at least one newly active endpoint constraint is in this list. All enumerations have fixed size; rational arithmetic has polynomial bit cost. This is not a uniform polynomial bound when $H$ or $n$ varies.
\end{proof}

\section{From finite local models to long improving paths}\label{sec:lifting}
We first isolate the exact obstruction to the proposed stronger certificate tests. This matrix decomposition is classical, not a new theorem of this paper.

\begin{lemma}[PSD principal block; Shaked-Monderer \cite{SM}, Lemma 3.1]\label{lem:spn}
If $Z\in\COP_n$ and $Z_{[r],[r]}\succeq0$, then $Z\in\SPN_n:=\mathbb S_+^n+\mathbb N^n$, where $\mathbb N^n$ denotes the symmetric entrywise nonnegative matrices. Hence $\ip A Z\ge0$ for every doubly nonnegative $A$.
\end{lemma}
\begin{proof}
First let $Z=\left(\begin{smallmatrix}G&b\\b\trans&c\end{smallmatrix}\right)$ with $G\succ0$. Minimize $x\trans Gx+2b\trans x$ on $x\ge0$, with minimizer $z$ and residual $\ell=Gz+b\ge0$, $z\trans\ell=0$. Copositivity at $(z,1)$ implies $c\ge z\trans Gz$. The decomposition
\[
 Z=\begin{pmatrix}G&-Gz\\-z\trans G&z\trans Gz\end{pmatrix}
   +\begin{pmatrix}0&\ell\\\ell\trans&c-z\trans Gz\end{pmatrix}
\]
is PSD plus nonnegative. For singular $G$, apply this result to $Z+\varepsilon\diag(I_r,0)$ and let $\varepsilon\downarrow0$. The SPN cone is closed: in a convergent sequence $Z_j=P_j+N_j$, diagonal entries of $P_j\succeq0$ and $N_j\ge0$ are bounded; $|(P_j)_{ab}|\le\sqrt{(P_j)_{aa}(P_j)_{bb}}$ bounds $P_j$, and then $N_j=Z_j-P_j$ is bounded. A convergent subsequence gives a limiting decomposition. The pairing assertion follows from PSD and entrywise nonnegative duality.
\end{proof}

\begin{proof}[Proof of Theorem~\ref{thm:main}]
Write $A=(B_N+Z)/(2\beta N^2)$, where $B_N=S_NBS_N$ and $\norm Z_F\le\eta$. Since $N\ge1$, $B_N\succeq\mu I$ and every entry is at least $b_*$. The choice of $\eta$ preserves positive definiteness and strict entrywise positivity. Also $\norm A_F\le(\beta N^2+1)/(2\beta N^2)\le1$.

Positive diagonal congruence preserves copositivity. Thus $W_N$ has cone margin $t_B/N^2$, norm at most $w_B$, and
\[
 \ip{W_N}{B_N}\le-\gamma+\frac{t_B}{N^2}\tr(B_N)
 \le-3\gamma/4.
\]
The perturbation contributes at most $\eta w_B\le\gamma/4$. Division by $2\beta N^2$ gives the stated pairing and distance. In particular every ball input is outside the ordinary CP cone.

Suppose a vertex is reachable in at most $m$ bounded edges within $\mathcal F_H$ from one of the allowed starts. Lemma~\ref{lem:gap} gives
\[
 \norm b_\infty\le R_0+mL_H,\qquad
 \norm k_\infty\le\norm z_\infty
 \le K_H(R_0+mL_H+1/2),
\]
so $v_k\le s_m$. For a primitive integral extreme ray $R$, define
\[
 h(R)=\ip{B_0}{R_{[r],[r]}}.
\]
All leading active constraints of $H$ remain active. Since $B_0$ is a positive combination of all their normals,
\begin{equation}\label{eq:exposure}
 h(R)\ge0,\qquad h(R)=0\ \Longleftrightarrow\ R_{[r],[r]}=0.
\end{equation}
The reverse implication in the equality uses that the active squares span $\mathbb S^r$. If $h(R)>0$, integrality gives $h(R)\ge1/q_B$. Applying the sharper block bounds \eqref{eq:raybounds},
\[
 \begin{split}
 \ip{B_N+Z}R
 &\ge\frac{N^2}{q_B}
       -2N\norm a_1F_Hs_m
       -(\alpha+\eta n)F_Hs_m^2\\
 &\ge\frac{N^2}{2q_B}>0.
 \end{split}
\]
For the last inequality, $N\ge8q_BF_HT_Bs_m$, $F_H,q_B\ge1$, $\eta\le1$, and $T_B\ge\norm a_1,\alpha+n,1$ bound the two subtracted fractions by $1/4$ and $1/64$ of $N^2/q_B$, respectively. Thus every improving extreme ray has zero leading block. The position-dependent quadratic growth of its last entry does not require a quadratic scale $N$: it is dominated by the $N^2$ leading exposure once $N$ is linear in the shear parameter.

There is an improving extreme ray at every departure vertex. Otherwise $A$ belongs to the minimizing normal cone, a finite cone generated by active nonnegative integer squares, and hence to $\CP_n$, a contradiction. The tangent cone is full dimensional and pointed by local polyhedrality and perfection. Its improving extreme ray has zero leading block and is therefore tangent to the face. It cannot be a recession direction: by Lemma~\ref{lem:gap}, those directions have nonnegative cross and bottom entries, so the strictly entrywise positive objective has positive cost on every nonzero such direction. A finite endpoint exists by closedness and local polyhedrality throughout the face; it is the adjacent vertex. Lemma~\ref{lem:spn} prevents a negative pairing at the current vertex, and also rules out every copositive certificate with PSD leading block. This proves the induction for all choices through $m$ pivots.

For the stated history certificate, its leading block is $(\sum_j\lambda_j)H+Q_{[r],[r]}\succeq0$, because every visited ray has zero leading block. Lemma~\ref{lem:spn} applies even if the procedure optimizes the correction or the history coefficients using unrestricted computation. Arbitrary nonnegative diagonal corrections are included; in fact arbitrary PSD corrections are included. The argument does not merely bound a particular correction formula.

All constants are fixed for a fixed block, seed, and starting region, and $N_m=\Theta(m+1)$. Entries of $A_{N_m}$ and $W_{N_m}$ have $O(\log(m+1))$ bits. The last diagonal of $A_{N_m}$ is a fixed positive rational multiple of $N_m^{-2}$, giving the matching lower encoding bound after reducing fractions. The distance estimate gives the same upper bound on logarithmic inverse distance. The local-cost claim is Corollary~\ref{cor:localcost}. Restarting within the starting region repeats an uncompleted prefix rather than shortening it; history certificates of the displayed form remain blocked across runs.
\end{proof}

\begin{remark}[An abstract exposure lemma]
The confinement argument is not peculiar to matrix cones. On a rational locally polyhedral set, let a face be exposed on its tangent cones by an integral functional $h$ with gap $1/q$ on primitive extreme rays outside the face. If local models yield transverse components bounded by $C(1+t)$ and $C(1+t)^2$ within graph distance $t$, the objective $N^2h+N\ell+a$ excludes every escaping improving extreme ray for $N\ge C'(1+t)$. The matrix theorem supplies verifiable hypotheses, the finite local models, and certificate exclusion for an important nonpolyhedral cone; it is not merely a formal restatement of this elementary domination argument.
\end{remark}

\section{Qualifying seeds without a five-cycle obstruction}\label{sec:seeds}
The criterion in Definition~\ref{def:seed} is open in the seed data. The following class demonstrates that its scope is not restricted to rescaling or embedding the original order-five example. For every odd order it produces a nonmember all of whose proper principal submatrices are CP, even after making the matrix dense.

Let $n\ge5$ be odd, let $G$ be the cycle on $(1,2,\ldots,n)$, and put
\begin{equation}\label{eq:cycle}
 t=\frac12+\frac1{16n^2},\quad F=I_n+t\operatorname{Adj}(G),\quad
 M=I_{n-1}-t\operatorname{Adj}(P_{n-1}),\quad h=M^{-1}\mathbf1.
\end{equation}
Here $P_{n-1}$ is the path on the first $n-1$ coordinates. Define
\[
 D=\diag(h_1,\ldots,h_{n-1},1),\quad
 C_{ij}=\min(i,j)(n+1-\max(i,j)),\quad
 \varepsilon=\frac1{32n^5},\quad B=DFD+\varepsilon C.
\]

\begin{theorem}[Arbitrarily large principal-minimal seeds]\label{thm:cycles}
The matrix $B$ is a qualifying seed for the canonical $H=T_{n-1}$. Every proper principal submatrix of $B$ belongs to the interior of its ordinary CP cone. It has an explicit rational weighted nonnegative Gram factorization. The normal-cone lifting theorem therefore gives canonical all-rule exponential lower bounds in every odd order, while all proper principal CP tests fail to detect nonmembership. By shrinking $\eta$ in \eqref{eq:robust} to at most $\varepsilon/8$, these properties persist throughout the entire robustness ball of the lifted targets.
\end{theorem}
\begin{proof}
The least eigenvalue of $F$ is $1-2t\cos(\pi/n)$. Using
$\cos(\pi/n)\le1-2/n^2$, obtained from the chord bound for $\sin$, gives
\[
 \lambda_{\min}(F)\ge\frac{15}{8n^2}+\frac1{4n^4}>\frac1{n^2}.
\]
The path matrix $M$ has the same least eigenvalue. Its inverse is nonnegative by the convergent Neumann series, and $h\ge\mathbf1$ is rational. In the leading block of $DFD$, the adjacent-pair interval square on $(i,i+1)$ has weight $t h_i h_{i+1}$. The remaining singleton coefficient is
\[
 h_i^2-t h_i\sum_{j:\,|i-j|=1}h_j=h_i>0.
\]
The canonical anchor $C$ is the sum of all consecutive-interval squares. Its leading block has coefficient at least one on every top-interval square. Thus $B_0\in\operatorname{int}\mathcal V_{T_{n-1}}$, and every entry of $B$ is positive.

For any triangle-free graph, $Q=J-2\operatorname{Adj}(G)$ is copositive. A self-contained Motzkin--Straus argument \cite{MS} is as follows: maximize the edge sum on the nonnegative simplex and choose a maximizer of minimum support. Nonadjacent positive coordinates can be merged into one of them without decreasing the sum, so the support is a clique. It has at most two vertices and the edge sum is at most $1/4$. Therefore $Q[x]\ge0$.
For the cycle,
\[
 \ip QF=n(1-2t)=-\frac1{8n}.
\]
Set $W=D^{-1}QD^{-1}$. Since $h_i\ge1$, $\norm W_F\le n$ and $\norm C_F\le\tr C\le n^3$. Consequently
\[
 \ip WB\le-\frac1{8n}+\varepsilon n^4=-\frac3{32n}.
\]
This proves nonmembership; one may use $\gamma=1/(16n)$, $\omega=n$, and $\mu=1/n^2$ as rational certificate parameters.

Let $S$ be a proper set of coordinates. The induced graph is a disjoint union of paths. Its comparison matrix
$M_S=I-t\operatorname{Adj}(G[S])$ is positive definite, since its path components have spectral radius at most $2\cos(\pi/n)$. Put $v=M_S^{-1}\mathbf1>0$. Then the following rational weighted factorization is exact:
\begin{equation}\label{eq:properfactor}
 F[S]=\sum_{i\in S}\frac1{v_i}e_ie_i\trans
 +\sum_{\{i,j\}\in E(G[S])}t v_i v_j
       (e_i/v_i+e_j/v_j)(e_i/v_i+e_j/v_j)\trans.
\end{equation}
Indeed the off-diagonal entries are $t$, and the diagonal at $i$ is
$(1+t\sum_{j\sim i}v_j)/v_i=1$. Congruence by $D[S]$ gives factors of $(DFD)[S]$. Adding the restricted interval-square factors of $\varepsilon C[S]$ gives factors of $B[S]$.

Moreover, these anchor atoms contain every consecutive-interval indicator in the induced ordering of $S$, with coefficient at least $\varepsilon$. Their squares span $\mathbb S^{|S|}$. Thus the factorization contains a positive reserve on a full rational basis and $B[S]\in\operatorname{int}\CP_{|S|}$. For a perturbation $Z$ with $\norm Z_F\le\varepsilon/8$, express $Z[S]$ in that interval basis. If entries outside the index range are set to zero, the coefficient of $[i,j]$ is
\[
 Z_{ij}-Z_{i-1,j}-Z_{i,j+1}+Z_{i-1,j+1}.
\]
Its magnitude is at most $4\norm Z_F\le\varepsilon/2$, so the positive reserve remains. For the lifted target use $S_N^{-1}ZS_N^{-1}$, whose norm cannot exceed $\norm Z_F$. Positive scaling and diagonal congruence preserve complete positivity and its interior. This proves the robust assertion for all proper subsets at once.
\end{proof}

The obstruction in order $n\ge7$ is therefore not a hidden order-five witness: every $5\times5$ principal submatrix, and even every $(n-1)\times(n-1)$ principal submatrix, is CP. Any copositive separator supported on a proper coordinate set has nonnegative pairing with the target. The difficulty is intrinsically global at that order. This remains a statement about \emph{ordinary} complete positivity, with no restrictions on Gram columns.

\begin{corollary}[Weighted cycle class]
Let $F$ be any rational positive-definite nonnegative matrix supported on an odd cycle of order $n\ge5$ and its diagonal. Suppose
\[
 \tr F<2\sum_{\{i,j\}\in E(G)}F_{ij}.
\]
There is a rational positive diagonal scaling of its first $n-1$ coordinates and an arbitrarily small positive rational canonical-anchor addition that produce a qualifying seed for $T_{n-1}$. All proper principal submatrices remain CP, and after the anchor addition they are interior CP. Thus the family includes a full range of nonuniform edge weights and diagonal data, not just the uniform matrices in \eqref{eq:cycle}.
\end{corollary}
\begin{proof}
Deleting the last vertex leaves a path. The leading comparison matrix, obtained by negating its off-diagonal entries, is diagonally sign-congruent to the leading block of $F$, hence positive definite. Its inverse is nonnegative: write it as $sI-N$ with $N\ge0$ and $s>\rho(N)$. Take $h$ to be its inverse applied to $\mathbf1$. The same scaled diagonal-dominance identity gives positive singleton and adjacent-pair interval coefficients. The witness $J-2\operatorname{Adj}(G)$ pairs negatively by the displayed condition. Preserve that strict gap with a small anchor addition. Every proper principal graph is a forest; the same comparison-matrix construction, with its original diagonal entries in place of one, gives a weighted factorization as in \eqref{eq:properfactor}. The anchor supplies interiority.
\end{proof}

For even order $n\ge6$, one may insert a positive diagonal padding block immediately before the distinguished last coordinate of an odd-order seed, and add a sufficiently small full canonical anchor. The old leading path occupies consecutive coordinates, and the distinguished old last coordinate remains last. The resulting leading block still has a positive interval representation and the embedded seed witness still separates. This recovers all fixed orders $n\ge5$. The stronger assertion that \emph{every} proper principal submatrix is CP is made for the odd-order construction, not for such padded examples.

\section{An independent theorem on intrinsic graph geometry}\label{sec:graph}
Let $\mathscr G_H$ denote the graph whose vertices are the vertices of $\mathcal E_H$ and whose edges are its bounded one-dimensional faces. Distances below are intrinsic to this graph. They need not equal distances in a larger Ryshkov graph that permits leaving the fixed-block face.

\begin{theorem}[Graph growth and lattice embedding]\label{thm:graph}
For every real $H\succ0$ of order $r$, the graph $\mathscr G_H$ is connected and has uniformly bounded degree. For every vertex $v$, there are positive constants $c_{H,v},C_{H,v}$ and $R_{H,v}$ such that
\[
 c_{H,v}R^r\le |B_{\mathscr G_H}(v,R)|\le C_{H,v}(R+1)^r
 \qquad(R\ge R_{H,v}).
\]
It contains a quasi-isometrically embedded copy of $\Z_+^r$: there are vertices $v_a$, $a\in\Z_+^r$, and constants $c_H,C_H>0$ with
\[
 c_H\norm{a-b}_\infty\le d_{\mathscr G_H}(v_a,v_b)
 \le C_H\norm{a-b}_1.
\]
When $H$ is rational, there are pairs of vertices of $O_H(\log(N+1))$ binary encoding length at distance $\Theta_H(N)$. Thus exponential graph distance in binary height is a geometric phenomenon independent of a separating objective or any five-cycle construction.
\end{theorem}
\begin{proof}
Uniform degree follows from Theorem~\ref{thm:templates}. The edge bound implies that a radius-$R$ graph ball about $(b_0,c_0)$ has
$\norm b_\infty\le\norm{b_0}_\infty+L_HR$.
At every vertex, $z=H^{-1}(\ell-b)$ and $\ell\in[0,1/2]^r$. Hence its shear index $k=\lfloor z\rfloor$ ranges over $O_{H,v}((R+1)^r)$ possibilities. There are only finitely many normalized templates, proving the upper bound.

We record connectivity carefully, since an infinite locally polyhedral set cannot be treated as a finite polytope without justification. Every graph vertex with $c\le C$ has $H[z]\le C-\kappa$, hence bounded $z$, and then bounded $b=\ell-Hz$. The finite-template theorem implies that there are only finitely many such vertices. Since the constraint $(u,y)=(0,1)$ gives $c\ge1$, and every $b\ge0$ gives $g_H(b)=1$, the minimum of $c$ is one. At a vertex with $c>1$ there is an extreme tangent direction decreasing $c$; otherwise the vertex minimizes $c$ globally. It is not a recession direction, so it gives a bounded edge. Strict descent through a finite set of vertices reaches the face $c=1$.

That face is a finite pointed polyhedron in the $b$ variables. Indeed, whenever $g_H(b)=1$, the gap estimate bounds $z$; every constraint active there has $y\le Y$ and $u=yz+e$ bounded. Thus only finitely many inequalities can become active anywhere on the face. These inequalities suffice globally: start at a strictly positive $b$, where every nonzero-$u$ constraint is strict, and take a segment to any alleged point satisfying the finite list but violating another constraint. At the first exit, local polyhedrality supplies an active defining constraint, which is one of the finite list, a contradiction. The recession cone of the face is $\R_+^r$. A positive linear objective has compact lower level sets on it, and after a generic perturbation uniquely minimizes at a vertex. Every other vertex has a bounded decreasing edge to that vertex, because the polyhedron has finitely many vertices and recession costs are positive. This proves connectivity of its vertex graph, and hence of $\mathscr G_H$.

For the lower bound we use an auxiliary unrestricted-integer epigraph
\[
 g_H^{\Z}(b)=\sup_{u\in\Z^r,\ y\ge1}\frac{1-H[u]-2yb\trans u}{y^2},
 \qquad \mathcal E_H^{\Z}=\{c\ge g_H^{\Z}(b)\}.
\]
Here $z=-H^{-1}b$, $\ell=0$, and the same gap, denominator and finite-template proofs apply, now with $k\in\Z^r$. Every $\mathsf S_k$ from \eqref{eq:shear} is an actual symmetry, since $u\mapsto u+yk$ permutes the integer numerators. The recession cone consists only of nonnegative vertical directions. The sublevel estimate
$g_H^{\Z}(b)\ge b\trans H^{-1}b+\kappa$ gives finitely many vertices at each bounded height. The minimum-$c$ face is $c=1$ and is bounded: the inequalities for $u=\pm e_i,y=1$ imply $|b_i|\le H_{ii}/2$. It is a nonempty finite polytope, by the same finite-active-list argument. Its graph is connected; descending $c$ proves that the entire bounded-edge graph is connected as well.

Choose a vertex $p$ of this auxiliary graph and, for each $i$, a finite edge path from $p$ to $\mathsf S_{e_i}p$. There are finitely many vertices in these paths and finitely many active vectors at each. Choose an integer $K\in\Z_+^r$ so large that $u+yK\ge0$ for all their active vectors. The same is true after adding any $a\in\Z_+^r$ to $K$. Each shifted vertex is a vertex of the orthant epigraph: it is feasible there, and all its spanning auxiliary active constraints are retained. Its orthant and unrestricted tangent cones coincide, because the auxiliary active set is fully retained. Consequently every shifted edge of each chosen path is also a bounded edge of $\mathcal E_H$. More explicitly, at the first endpoint the direction remains an extreme ray, and the same next endpoint retains the active constraint that prevents extension.

Define $v_a=\mathsf S_{K+a}p$. Let $D$ be the largest length of the finitely many chosen generator paths. Translating these paths and reversing them when necessary gives
\[
 d_{\mathscr G_H}(v_a,v_b)\le D\norm{a-b}_1.
\]
One may change coordinates one at a time through nonnegative lattice indices, so no translated path requires a negative index. Their cross coordinates differ by $H(a-b)$. Along any graph path, \eqref{eq:movement} gives
\[
 d_{\mathscr G_H}(v_a,v_b)
 \ge\frac{\norm{H(a-b)}_\infty}{L_H}
 \ge\frac{\norm{a-b}_\infty}{K_HL_H}.
\]
This proves the embedding and supplies $\Omega_H(R^r)$ vertices within radius $R$ of $v_0$. Connectivity transfers the growth bound to any root with an additive distance shift.

For rational $H$, the auxiliary vertex $p$ is rational, being determined by rational active equations. Its shear by $K+Ne_1$ has coordinates polynomial of degree at most two in $N$ with fixed rational coefficients. Its encoding is $O_H(\log(N+1))$, while the two distance bounds give $\Theta_H(N)$. This proves the last assertion.
\end{proof}

\begin{corollary}[A complete elementary model]\label{cor:grid}
For $H=I_2$, the graph $\mathscr G_H$ is exactly the nearest-neighbor graph of $\Z_+^2$. Its vertices are
\[
 b_i=-k_i-\tfrac12,\qquad c=1+k_1(k_1+1)+k_2(k_2+1),
 \qquad(k_1,k_2)\in\Z_+^2.
\]
The graph distance is $\ell^1$ distance and the ball of radius $R$ about $(k_1,k_2)=(0,0)$ has $\binom{\lfloor R\rfloor+2}{2}$ vertices.
\end{corollary}
\begin{proof}
Let $z_i=\max(0,-b_i)$. Rounding each positive $z_i$ to its nearest nonnegative integer gives a denominator-one value at least $\norm z_2^2+1/2$. Every denominator $y\ge2$ gives at most $\norm z_2^2+1/4$ and is therefore inactive. Thus
\[
 g_{I_2}(b)=1+\sum_{i=1}^2\max_{u_i\in\Z_+}(-u_i^2-2b_i u_i).
\]
A vertex requires a tie in each coordinate, since otherwise the active affine normals fail to span $\R^3$. Ties occur exactly between $k_i$ and $k_i+1$ at $b_i=-k_i-1/2$. To obtain a bounded edge one tie is retained and the other moves between successive breakpoints. There are no diagonal edges: changing both ties leaves only one common active slope, which is insufficient for a one-dimensional face in three dimensions. The outward edges at $k_i=0$ are unbounded and are excluded from the graph. The distance and ball count follow.
\end{proof}

The auxiliary unrestricted epigraph is related to classical inhomogeneous quadratic forms and Delaunay geometry \cite{Erdahl,IE}. The assertions here include the positive-orthant truncation, all integer last-coordinate denominators, bounded local models at boundary vertices, and an intrinsic graph-growth conclusion. Periodicity of an unrestricted denominator-one lattice diagram alone would not prove these assertions.

\section{Necessary restrictions, verification, and comparison}\label{sec:scope}
\subsection{What the assumptions do and do not assert}
The strict normal-cone condition is sufficient for the lifting theorem; we do not claim it is necessary for every form of local-search hardness. It is essential to this confinement mechanism. If $B_0\notin\mathcal V_H$, separation gives a leading tangent direction $X$ with $\ip{B_0}X<0$. At any fixed vertex of $\mathcal F_H$, extend it to a full tangent direction by choosing its bottom entry sufficiently large: all finitely many positive-last-coordinate active inequalities can then be satisfied, while the leading active inequalities already hold. Decomposing the pointed tangent cone into its extreme rays, at least one ray has negative leading cost. For large $N$, that fixed ray improves the scaled objective and leaves the face.

Boundary exposure can also fail. At $T_3$, the direction
\[
 R_{12}=R_{21}=1/2,\qquad R_{13}=R_{31}=-1/2,
\]
with all other entries zero is an extreme tangent ray: among the six interval inequalities only the interval $[1,2]$ has positive value. With $B=I_3+(e_1e_3\trans+e_3e_1\trans)/4$, the leading block $I_2$ lies on the boundary of $\mathcal V_{T_2}$ and $\ip{S_NBS_N}R=-N/4<0$. This example diagnoses the exposure condition; it is not claimed to be a nonmember lower-bound example.

Multiple starts sharing the leading block are covered in full. Arbitrary menus of distinct leading perfect blocks are not covered. If $H_1\ne H_2$ are distinct perfect vertices, the interiors of their minimizing normal cones are disjoint: a functional in either interior uniquely minimizes at its vertex. Thus one leading seed cannot satisfy both strict exposure conditions. This is an obstruction to this proof strategy, not a theorem that all diverse-start methods are efficient or that no other hardness construction exists.

The extreme-ray restriction is also real. Adding a sufficiently small positive multiple of an escaping feasible tangent direction to a strictly improving face ray still gives a feasible improving direction, but leaves the face. Consequently this proof does not cover arbitrary nonextreme descent directions, nonlocal jumps, or input-dependent initializations. Arbitrary global certificate generation is intentionally outside the lower bound; $W_N$ is a short explicit global separator.

Positive definiteness is needed for the independent geometric theorem. For example $H=0$ gives $\mathcal E_H=\{b\ge0,c\ge1\}$, which has only one vertex rather than polynomial graph growth of positive degree. All algorithmic claims with fixed $H,n$ allow constants depending on both; the template bounds can be very large and do not imply uniformly practical computation in variable dimension.

\subsection{Relationship to prior work}
The simplex model, perfect matrices, finite normal cones, and qualitative Ryshkov geometry are antecedents \cite{DSV,DS,OS}. The bounded-edge displacement estimate and the canonical five-cycle all-rule construction were proved in the preceding project manuscript, and are reproved here rather than relabeled as new. The strengthening is the finite-template normalization and its sharper ray bounds, the seed-based linear-scale lifting theorem, its stronger certificate exclusions, higher-order principal-minimal applications, and the independent graph-growth theorem.

The PSD-principal-block lemma is explicitly attributed to \cite{SM}; it is not a new matrix decomposition. Odd-cycle obstructions and the distinction between doubly nonnegative and completely positive matrices are classical. The new odd-cycle conclusion couples these seeds to all-choice canonical hardness and retains complete positivity of every proper principal block, even on a specified neighborhood. Motzkin--Straus supplies the classical graph witness; a self-contained proof was given above. The unrestricted lattice geometry is related to \cite{Erdahl,IE}; our conclusions concern a different, orthant-truncated rational epigraph and its full adjacency graph. The recent algorithmic work \cite{Minimum} studies the copositive minimum rather than a lower bound on adjacent traversal.

These comparisons isolate the proposed new assertions. They are not an exhaustive first-priority certificate. A lack of an identical result in the sources inspected is not a proof of historical novelty. The earlier suitable-pivot finite-termination conjecture is not contradicted: an exponentially long path may still terminate. We do not assert eventual termination of every permitted local rule.

\subsection{Reproducibility and status}
The supplement contains exact rational seed certificates, finite normal-system checks, shear-template checks, proper-principal CP factors, and the two-dimensional graph model. A separately written verifier imports no generator code and rejects deliberately corrupted certificates. The original canonical certificate suite is retained and rerun without modifying its formulas. The verification report records the actual case counts and file hashes; finite tests supplement, rather than replace, the proofs.

No production high-dimensional polyhedral library, external human proof review, or proof-assistant formalization is claimed. The user-reported randomized ten-paper benchmark is not recalibrated in this mathematical supplement: the sample and its complete comparison rubric were not available in the retrieved research files. The revised contribution is suitable for a fresh skeptical comparison, but the paper does not assert a top-three rank or an editorial probability.

\begin{thebibliography}{9}
\bibitem{DSV} M. Dutour Sikiri\'c, A. Sch\"urmann, F. Vallentin, A simplex algorithm for rational cp-factorization, Math. Program. 187 (2021), 25--45. doi:10.1007/s10107-020-01467-4.
\bibitem{DS} V. Dannenberg, A. Sch\"urmann, Perfect copositive matrices, Commun. Math. 31 (2023), 137--156. doi:10.46298/cm.11141.
\bibitem{OS} A. Oertel, A. Sch\"urmann, Generalized Perfect Matrices, arXiv:2602.05841v2 (2026).
\bibitem{SM} N. Shaked-Monderer, SPN graphs: when copositive = SPN, Linear Algebra Appl. 509 (2016), 82--113. doi:10.1016/j.laa.2016.07.018. Lemma 3.1; arXiv:1604.02172v3.
\bibitem{MS} T. S. Motzkin, E. G. Straus, Maxima for graphs and a new proof of a theorem of Tur\'an, Canad. J. Math. 17 (1965), 533--540. doi:10.4153/CJM-1965-053-6.
\bibitem{Erdahl} R. M. Erdahl, A cone of inhomogeneous second-order polynomials, Discrete Comput. Geom. 8 (1992), 387--416. doi:10.1007/BF02293055.
\bibitem{IE} M. Dutour Sikiri\'c, A. Sch\"urmann, F. Vallentin, Inhomogeneous extreme forms, Ann. Inst. Fourier 62 (2012), 2227--2255. doi:10.5802/aif.2748.
\bibitem{Minimum} A. Oertel, A. Sch\"urmann, On Computing the Copositive Minimum and its Representatives, arXiv:2509.23696 (2025).
\end{thebibliography}
\end{document}
